\documentclass[11pt,a4paper]{article}

% ─── Paquetes ──────────────────────────────────────────────
\usepackage[utf8]{inputenc}
\usepackage[T1]{fontenc}
\usepackage{lmodern}
\usepackage[margin=2.3cm]{geometry}
\usepackage{amsmath,amssymb}
\usepackage{graphicx}
\usepackage{xcolor}
\usepackage{hyperref}
\usepackage{booktabs}
\usepackage{array}
\usepackage{tabularx}
\usepackage{enumitem}
\usepackage{fancyhdr}
\usepackage{titlesec}
\usepackage{listings}
\usepackage{tcolorbox}
\usepackage{fontawesome5}
\usepackage{multicol}
\usepackage{parskip}

% ─── Colors ────────────────────────────────────────────────
\definecolor{accent}{HTML}{B07A00}
\definecolor{greenhl}{HTML}{1A7F37}
\definecolor{bluehl}{HTML}{0969DA}
\definecolor{redhl}{HTML}{CF2626}
\definecolor{warnhl}{HTML}{9A6700}
\definecolor{lightbg}{HTML}{FDF6E3}
\definecolor{surf2}{HTML}{F5EDCE}
\definecolor{border}{HTML}{DDD0A8}

% ─── Configuración de Hyperref ──────────────────────────────
\hypersetup{
  colorlinks=true,
  linkcolor=accent,
  urlcolor=bluehl,
  citecolor=greenhl,
  pdftitle={Weibull-Hyperbola Modelo v8 - Manual de usuario},
  pdfauthor={},
  pdfsubject={Manual de usuario}
}

% ─── Encabezado / pie de página ────────────────────────────
\pagestyle{fancy}
\fancyhf{}
\fancyhead[L]{\small\textcolor{accent}{\textbf{Weibull--Hyperbola Modelo}}}
\fancyhead[R]{\small v8}
\fancyfoot[C]{\small\thepage}
\renewcommand{\headrulewidth}{0.4pt}

% ─── Sección styling ────────────────────────────────────────
\titleformat{\section}
  {\Large\bfseries\color{accent}}{\thesection}{0.6em}{}
\titleformat{\subsection}
  {\large\bfseries\color{black!80}}{\thesubsection}{0.5em}{}
\titleformat{\subsubsection}
  {\normalsize\bfseries\color{black!70}}{\thesubsubsection}{0.4em}{}

% ─── Recuadros personalizados ──────────────────────────────
\newtcolorbox{notebox}[1][]{
  colback=lightbg, colframe=accent, boxrule=0.6pt,
  arc=2mm, left=4mm, right=4mm, top=2mm, bottom=2mm,
  fonttitle=\bfseries\small, title={#1}
}
\newtcolorbox{warnbox}[1][]{
  colback=orange!6, colframe=warnhl, boxrule=0.6pt,
  arc=2mm, left=4mm, right=4mm, top=2mm, bottom=2mm,
  fonttitle=\bfseries\small, title={#1}
}
\newtcolorbox{tipbox}[1][]{
  colback=green!5, colframe=greenhl, boxrule=0.6pt,
  arc=2mm, left=4mm, right=4mm, top=2mm, bottom=2mm,
  fonttitle=\bfseries\small, title={#1}
}

% ─── Listados de código ────────────────────────────────────
\lstset{
  basicstyle=\ttfamily\small,
  backgroundcolor=\color{lightbg},
  frame=single,
  rulecolor=\color{border},
  breaklines=true,
  columns=fullflexible,
  keepspaces=true,
}

% ─── Macros auxiliares ─────────────────────────────────────
\newcommand{\piz}{\ensuremath{\pi_3}}
\newcommand{\pit}{\ensuremath{\pi_2}}
\newcommand{\pio}{\ensuremath{\pi_1}}
\newcommand{\sm}{\ensuremath{\sigma_m}}
\newcommand{\sM}{\ensuremath{\sigma_M}}
\newcommand{\sMzero}{\ensuremath{\sigma_{M0}}}
\newcommand{\Nzero}{\ensuremath{N_0}}
\newcommand{\Weibull}{\textsc{Weibull}}

\setlength{\parindent}{0pt}
\setlength{\parskip}{5pt}

\begin{document}

% ════════════════════════════════════════════════════════════
% TITLE PAGE
% ════════════════════════════════════════════════════════════
\begin{titlepage}
\centering
\vspace*{3cm}

{\Huge\bfseries\color{accent} Weibull--Hyperbola Modelo}\\[0.4cm]
{\LARGE\bfseries Manual de usuario}\\[1.5cm]

{\large Version 8}\\[0.3cm]
{\large Aplicación web interactiva para el análisis de datos de fatiga}\\[3cm]

\begin{tabular}{rl}
\textbf{Tipo de aplicación:} & Aplicación web de una sola página \\
\textbf{Modeloos:} & M0--M7 (8 modelos anidados) \\
\textbf{Estimación:} & Máxima verosimilitud + búsqueda en rejilla \\
\textbf{Simulación:} & Rama física garantizada \\
\textbf{Exportaraciones:} & CSV, LaTeX, PDF, PNG \\
\end{tabular}

\vfill

{\small Este documento describe cómo utilizar la aplicación web del Modeloo Weibull--Hipérbola para ajustar modelos, simular y analizar datos de vida a fatiga.}

\end{titlepage}

% ════════════════════════════════════════════════════════════
\tableofcontents
\newpage

% ════════════════════════════════════════════════════════════
\section{Introducción}
% ════════════════════════════════════════════════════════════

El \textbf{Modelo Weibull--Hipérbola} es una herramienta basada en navegador para analizar datos de fatiga. Combina una relación hiperbólica tensión--vida con una distribución \Weibull{} de tres parámetros para modelar la dispersión de la vida a fatiga.

La aplicación permite:
\begin{itemize}[leftmargin=1.4em]
  \item Simular datos sintéticos de fatiga en una rama física garantizada.
  \item Ajustar hasta ocho modelos anidados (M0--M7) a datos experimentales.
  \item Estimar automáticamente los parámetros de referencia $\sMzero$ y $\Nzero$.
  \item Comparar modelos mediante AIC y BIC.
  \item Visualizar curvas de percentiles tanto en coordenadas transformadas como originales.
  \item Exportarar resultados como CSV, tablas LaTeX y figuras PDF y PNG.
\end{itemize}

\begin{notebox}[No se requiere instalación]
La aplicación consta de un único archivo HTML. Se ejecuta íntegramente en el navegador: no necesita servidor ni instalación, y ningún dato sale de su ordenador. Todos los cálculos se realizan localmente en JavaScript.
\end{notebox}

% ════════════════════════════════════════════════════════════
\section{El modelo matemático}
% ════════════════════════════════════════════════════════════

\subsection{Variables transformadas}

El modelo utiliza tres variables adimensionales:
\begin{align}
  \pio &= \frac{\sm}{\sM} && \text{(razón de tensiones)} \\[2pt]
  \pit &= \frac{\sM}{\sMzero} - 1 && \text{(tensión de amplitud normalizada)} \\[2pt]
  \piz &= \log_{10}\!\left(\frac{N}{\Nzero}\right) && \text{(logaritmo de la vida normalizado)}
\end{align}
donde $\sm$ es la tensión media, $\sM$ es la amplitud de tensión, $N$ es el número de ciclos hasta el fallo, y $\sMzero$ y $\Nzero$ son valores de referencia.

\subsection{La hipérbola}

La relación fundamental es:
\begin{equation}
  V = w_1\pio + w_{12}\pio\pit + w_{13}\pio\piz + \pit\piz + w_{123}\pio\pit\piz
  \label{eq:V}
\end{equation}
Solving for $\pit$ as a function of $\piz$ at fixed $\pio$ gives a hyperbola:
\begin{equation}
  \pit = \frac{V - w_1\pio - w_{13}\pio\piz}{\piz + w_{12}\pio + w_{123}\pio\piz}
\end{equation}
La asíntota vertical está en $\piz^{\text{asym}} = -w_{12}\pio/(1+w_{123}\pio)$ y la asíntota horizontal en $\pit^{\text{asym}} = -w_{13}\pio/(1+w_{123}\pio)$.

\subsection{La distribución de Weibull}

La variable $V$ sigue una distribución \Weibull{} de tres parámetros:
\begin{equation}
  f(v) = \frac{\beta}{\delta}\left(\frac{v-\lambda}{\delta}\right)^{\beta-1}
         \exp\!\left[-\left(\frac{v-\lambda}{\delta}\right)^{\beta}\right], \quad v > \lambda
\end{equation}
con parámetro de posición $\lambda$, escala $\delta$ y forma $\beta$.

\subsection{Los ocho modelos anidados}

\begin{center}
\begin{tabular}{cl}
\toprule
\textbf{Modelo} & \textbf{Parámetros libres} \\
\midrule
M0 & $w_1, w_{12}, w_{13}, w_{123}$ \\
M1 & $w_1, w_{12}, w_{13}$ \\
M2 & $w_1, w_{13}, w_{123}$ \\
M3 & $w_1, w_{12}, w_{123}$ \\
M4 & $w_1, w_{123}$ \\
M5 & $w_1, w_{13}$ \\
M6 & $w_1, w_{12}$ \\
M7 & $w_1$ \\
\bottomrule
\end{tabular}
\end{center}

Los parámetros que no aparecen en la tabla se fijan en cero. Todos los modelos están anidados en M0.

% ════════════════════════════════════════════════════════════
\section{Primeros pasos}
% ════════════════════════════════════════════════════════════

\subsection{Abrir la aplicación}

Abra el archivo HTML en cualquier navegador moderno (Chrome, Firefox, Edge o Safari). Después de la primera carga no se necesita instalación ni conexión a Internet, salvo que las bibliotecas Chart.js y jsPDF se carguen desde archivos locales.

\subsection{Descripción de la interfaz}

La interfaz está organizada en secciones desplegables:

\begin{center}
\begin{tabularx}{\textwidth}{lX}
\toprule
\textbf{Sección} & \textbf{Función} \\
\midrule
\textcolor{bluehl}{\textbf{Modelo y parámetros}} & Elegir el modelo y editar los parámetros $w$, los parámetros de \Weibull{} y los valores de referencia. \\
\textcolor{greenhl}{\textbf{Configuración de simulación}} & Definir los intervalos de simulación y ejecutarla. \\
\textcolor{orange}{\textbf{Extremos de los datos}} & Visualización automática, solo de lectura, de los intervalos de los datos. \\
\textcolor{redhl}{\textbf{Fitting \& Estimation}} & Cargar un CSV, estimar referencias y ajustar modelos. \\
\textcolor{accent}{\textbf{Resultados}} & Pestañalas, gráficos y opciones de exportación. \\
\textcolor{purple}{\textbf{Registro de consola}} & Mensajes, advertencias e información sobre el progreso. \\
\bottomrule
\end{tabularx}
\end{center}

Haga clic en el encabezado de cualquier sección para desplegarla o contraerla. El botón superior derecho \textbf{Expandir todo / Contraer todo} cambia todas las secciones a la vez.

\subsection{El indicador de estado}

Un punto de color en la esquina superior derecha del encabezado indica el estado de la aplicación:
\begin{itemize}[leftmargin=1.4em]
  \item \textcolor{greenhl}{\textbf{Verde}} --- lista / inactiva.
  \item \textcolor{redhl}{\textbf{Rojo}} --- calculando. No cierre la pestaña.
\end{itemize}

\subsection{El campo Caso / Título}

El campo de texto del encabezado se utiliza como \textbf{prefijo del nombre de archivo} para todas las exportaciones. Al cargar un CSV, se establece automáticamente a partir del nombre del archivo.

% ════════════════════════════════════════════════════════════
\section{Flujo de trabajo 1: simulación de datos}
% ════════════════════════════════════════════════════════════

Este flujo de trabajo genera datos sintéticos a partir de parámetros conocidos.

\subsection{Paso 1 --- Seleccionar el modelo y los parámetros}

\begin{enumerate}[leftmargin=1.4em]
  \item Expand \textbf{Modelo \& Parameters}.
  \item Elija un modelo en el selector \textit{Modelo} (por ejemplo, M0). Los campos de los parámetros $w$ libres se actualizan automáticamente.
  \item Introduzca los valores de $w$. Un conjunto inicial razonable es $w_1=0.5$, $w_{12}=-0.2$, $w_{13}=0.3$, $w_{123}=-0.1$.
  \item Introduzca los parámetros de \Weibull{}: $\lambda=0.5$, $\delta=0.3$, $\beta=2.5$ (valores predeterminados).
  \item Introduzca los valores de referencia: $\sMzero=40$, $\Nzero=1000$ (defaults).
\end{enumerate}

\begin{warnbox}[Restricciones de los valores de referencia]
$\sMzero$ debe ser estrictamente menor que $\sM^{\min}$ en el intervalo de simulación. $\Nzero$ debe ser menor que el umbral de parada sin fallo. La aplicación corrige automáticamente los incumplimientos y los comunica en la consola.
\end{warnbox}

\subsection{Paso 2 --- Establecer los intervalos de simulación}

Despliegue \textbf{Configuración de simulación} y establezca:
\begin{itemize}[leftmargin=1.4em]
  \item $\sM^{\min}$, $\sM^{\max}$ --- intervalo de amplitud de tensión (e.g., 50--200 MPa).
  \item $\sm^{\min}$, $\sm^{\max}$ --- intervalo de tensión media (e.g., 10--150 MPa).
  \item \textbf{N points} --- número de puntos que se generarán (valor predeterminado: 1000).
  \item \textbf{Runout threshold} --- número de ciclos por encima del cual un punto se considera censurado (valor predeterminado: $10^7$).
  \item \textbf{Apply runout censoring} --- si se marca, los puntos que superen el umbral se etiquetan como ensayos sin fallo.
\end{itemize}

\subsection{Paso 3 --- Simular}

Haga clic en \textbf{\textcolor{greenhl}{Simular datos}}. La consola informa de:
\begin{itemize}[leftmargin=1.4em]
  \item Número de puntos válidos generados.
  \item Tasa de rechazo (porcentaje de intentos descartados por las restricciones físicas).
\end{itemize}

\begin{tipbox}[Interpretación de la tasa de rechazo]
Una tasa de rechazo inferior al 80\% es normal. Una tasa superior al 80\% sugiere que los parámetros $w$, los intervalos o los valores de referencia no son adecuados. La consola mostrará una advertencia.
\end{tipbox}

\subsection{Paso 4 --- Representar las curvas}

Haga clic en \textbf{\textcolor{accent}{Representar curvas (sim.)}} para mostrar las curvas de percentiles. Se utilizan los parámetros actuales de la sección \textbf{Modelo y parámetros}; no se realiza ningún ajuste.

% ════════════════════════════════════════════════════════════
\section{Flujo de trabajo 2: ajuste de datos experimentales}
% ════════════════════════════════════════════════════════════

Este flujo de trabajo estima parámetros a partir de datos de fatiga medidos.

\subsection{Paso 1 --- Preparar el archivo CSV}

El archivo debe estar separado por comas y contener las siguientes columnas obligatorias:

\begin{center}
\begin{tabular}{lll}
\toprule
\textbf{Columna} & \textbf{Descripción} & \textbf{Unidades} \\
\midrule
\texttt{sigmam} & Tensión media $\sm$ & MPa \\
\texttt{sigmaM} & Amplitud de tensión $\sM$ & MPa \\
\texttt{N} & Ciclos hasta el fallo & cycles \\
\bottomrule
\end{tabular}
\end{center}

Puede incluirse una columna opcional \texttt{runout} (0/1). Si no existe, se infiere a partir del umbral de parada sin fallo.

Ejemplo:
\begin{lstlisting}
sigmam,sigmaM,N
20,180,15000
30,160,42000
40,140,98000
...
\end{lstlisting}

\subsection{Paso 2 --- Cargar el archivo}

Puede hacerlo de una de estas dos formas:
\begin{itemize}[leftmargin=1.4em]
  \item Arrastre el archivo a la zona de carga, o
  \item Haga clic en la zona de carga para buscar el archivo.
\end{itemize}

Al cargarlo, la aplicación:
\begin{enumerate}[leftmargin=1.4em]
  \item Borra todos los resultados anteriores.
  \item Establece el selector \textit{Modelo} en \textbf{Todos} y marca las casillas de todos los modelos.
  \item Actualiza la sección \textbf{Extremos de los datos} con los mínimos y máximos de cada variable.
  \item Establece el \textbf{Caso / Título} a partir del nombre del archivo.
  \item Muestra la tabla de datos en la pestaña \textbf{Resultados}.
\end{enumerate}

\subsection{Paso 3 --- Configurar el ajuste}

Expand \textbf{Fitting \& Estimation}:

\begin{enumerate}[leftmargin=1.4em]
  \item \textbf{Modeloos que se ajustarán} --- marque los modelos que desea incluir (por defecto están todos marcados).
  \item \textbf{Método de estimación}:
    \begin{itemize}
      \item \textbf{Opción 1 (recomendada)} --- cada modelo estima de forma independiente $\sMzero$ y $\Nzero$.
      \item \textbf{Opción 2} --- estimación conjunta: un único par $\sMzero/\Nzero$ para todos los modelos y, después, se vuelven a ajustar.
    \end{itemize}
  \item \textbf{Configuración de la búsqueda en rejilla}:
    \begin{itemize}
      \item \textbf{Grid points} (3--20) --- resolución de la rejilla inicial gruesa.
      \item \textbf{Grid restarts} (1--15) --- reinicios de L-BFGS durante la fase de rejilla.
      \item \textbf{Refinement restarts} (3--30) --- reinicios de L-BFGS durante el refinamiento final.
    \end{itemize}
\end{enumerate}

\begin{notebox}[Consejo sobre rendimiento]
Los valores predeterminados (3, 3, 5) son rápidos. Para ajustes destinados a publicaciones, utilice entre 5 y 7 puntos de rejilla y entre 5 y 10 reinicios. Los valores mayores requieren proporcionalmente más tiempo.
\end{notebox}

\subsection{Paso 4 --- Estimar o ajustar}

Hay dos botones disponibles:

\begin{itemize}[leftmargin=1.4em]
  \item \textbf{\textcolor{warnhl}{\textbf{[Plot]} Estime $\sMzero$ \& $\Nzero$}} --- ejecuta la búsqueda completa en rejilla y el refinamiento.
  \item \textbf{\textcolor{accent}{Ajuste manual ($\sMzero$, $\Nzero$)}} --- ajusta los modelos con los valores de referencia introducidos en \textbf{Modelo y parámetros}, sin realizar la búsqueda.
\end{itemize}

Durante la ejecución, el punto de estado se vuelve \textcolor{redhl}{rojo}. El progreso se muestra en la barra de progreso y en la consola.

\subsection{Paso 5 --- Revisar los resultados}

Cuando termina el ajuste, la sección \textbf{Resultados} muestra cinco pestañas:

\begin{center}
\begin{tabularx}{\textwidth}{lX}
\toprule
\textbf{Pestaña} & \textbf{Contenido} \\
\midrule
\textbf{Datos} & Datos experimentales cargados (paginados). \\
\textbf{Modelo Comparison} & Estimaciones de parámetros, AIC y BIC para todos los modelos ajustados. \\
\textbf{Curvas de percentiles} & Interactive charts (see Sección~\ref{sec:charts}). \\
\textbf{Datos simulados} & Pestañala de datos simulados, si los hay. \\
\textbf{Exportar} & Botones para descargar CSV, LaTeX, PDF y PNG. \\
\bottomrule
\end{tabularx}
\end{center}

El mejor modelo (el que tiene el AIC más bajo) se resalta con fondo verde y un asterisco en la tabla de comparación. Su ecuación aparece encima de las pestañas.

% ════════════════════════════════════════════════════════════
\section{Las pestañas de resultados en detalle}
% ════════════════════════════════════════════════════════════

\subsection{Pestaña Datos}

Muestra los datos experimentales cargados con una columna adicional $R = \sm/\sM$ (razón de tensiones). Las tablas se presentan en páginas de 10 filas.

\subsection{Pestaña Comparación de modelos}

Muestra una fila por cada modelo ajustado, con las siguientes columnas:
\begin{center}
\begin{tabular}{ll}
\toprule
\textbf{Columna} & \textbf{Significado} \\
\midrule
\textbf{*} & Identifica el mejor modelo (AIC más bajo) \\
Modelo & Identificador del modelo M0--M7 \\
$k$ & Número de parámetros libres \\
$\sMzero$, $\Nzero$ & Valores de referencia estimados \\
$w_1,w_{12},w_{13},w_{123}$ & Parámetros $w$ estimados \\
$\lambda$, $\delta$, $\beta$ & \Weibull{} parameters \\
$\ell$ & Log-verosimilitud \\
AIC, BIC & Criterios de información \\
\bottomrule
\end{tabular}
\end{center}

\subsection{Pestaña Curvas de percentiles}
\label{sec:charts}

Muestra el modelo ajustado como una familia de curvas de percentiles ($p = 0.01, 0.05, 0.25, 0.50, 0.75, 0.95, 0.99$) at six values of $\pio$.

\subsubsection*{Sistemas de coordenadas}

Utilice el selector de \textbf{Modelo y parámetros} para alternar entre:
\begin{itemize}[leftmargin=1.4em]
  \item \textbf{$\pit$ vs $\piz$} --- coordenadas transformadas.
  \item \textbf{$\sM$ vs $N$} --- coordenadas físicas originales (valor predeterminado).
\end{itemize}

\subsubsection*{Control deslizante de opacidad}

El control deslizante situado encima de los gráficos determina cómo se atenúan los puntos de dispersión con la distancia al valor de la sección $\pio$:
\begin{itemize}[leftmargin=1.4em]
  \item \textbf{0} --- todos los puntos tienen la misma opacidad.
  \item \textbf{4.6} --- los puntos a distancia mediana son casi invisibles; los cercanos a la sección permanecen destacados.
\end{itemize}
Esto ayuda a visualizar qué puntos de datos pertenecen a cada sección de $\pio$.

\subsubsection*{Desplazamiento}

Los gráficos están dispuestos horizontalmente. Desplácese lateralmente (o deslice el dedo en dispositivos táctiles) para ver todas las secciones.

\subsubsection*{Datos superpuestos}

When \textbf{Incluir datos simulados} is active (see Sección~\ref{sec:include-sim}), the simulated points are shown in blue and the real experimental data in red.

\subsection{Pestaña Datos simulados}

Muestra los datos sintéticos generados, si se ha ejecutado una simulación. Se pagina igual que las demás tablas.

\subsection{Pestaña Exportarar}

Incluye:
\begin{itemize}[leftmargin=1.4em]
  \item \textbf{\textbf{[Download]} Datos CSV} --- datos experimentales en formato CSV.
  \item \textbf{\textbf{[Download]} Modelos CSV} --- estimaciones de parámetros en formato CSV.
  \item \textbf{\textbf{[Download]} Datos LaTeX} --- datos experimentales en una tabla con formato booktabs.
  \item \textbf{\textbf{[Download]} Modelos LaTeX} --- tabla de parámetros con formato booktabs.
  \item \textbf{\textbf{[Download]} Simulated CSV} --- datos simulados en formato CSV.
  \item \textbf{\textbf{[Plot]} Exportar Figures (PDF)} --- un PDF por gráfico, en formato A4 apaisado.
  \item \textbf{\textbf{[Plot]} Exportar Figures (PNG)} --- un PNG por gráfico, con resolución 3×.
\end{itemize}

Todos los nombres de archivo utilizan \textbf{Caso / Título} como prefijo. Debajo de los botones aparecen miniaturas de las figuras; al hacer clic en una miniatura, se descarga la figura como PNG.

% ════════════════════════════════════════════════════════════
\section{Parámetros de referencia y autocorrección}
% ════════════════════════════════════════════════════════════

\subsection{Restricciones}

\begin{itemize}[leftmargin=1.4em]
  \item $\sMzero$ debe ser \textbf{estrictamente menor} que el mínimo de $\sM$ en los datos (de lo contrario, $\pit \le 0$).
  \item $\Nzero$ debe ser \textbf{estrictamente menor} que el mínimo de $N$ en los datos (de lo contrario, $\piz \le 0$).
\end{itemize}

\subsection{Autocorrección}

Al hacer clic en \textbf{Simular datos} o \textbf{Ajuste manual}, si alguno de los valores de referencia incumple su restricción, la aplicación lo corrige automáticamente a $0.9\times$ el mínimo correspondiente y escribe una advertencia en la consola.

Al hacer clic en \textbf{Estimar $\sMzero$ y $\Nzero$}, los intervalos de búsqueda se eligen automáticamente dentro de la región factible.

\begin{notebox}[¿Por qué es importante?]
La hipérbola tiene dos ramas. Solo una proporciona datos de fatiga físicamente válidos ($\sM$ positiva y $N$ positivo). La aplicación impone la restricción $\piz + w_{12}\pio + w_{123}\pio\piz > 0$ para todos los puntos de datos, lo que garantiza la rama correcta.
\end{notebox}

% ════════════════════════════════════════════════════════════
\section{Incluir datos simulados en las curvas}
\label{sec:include-sim}
% ════════════════════════════════════════════════════════════

La opción \textbf{Incluir datos simulados} (en Configuración de simulación) tiene tres efectos:

\begin{enumerate}[leftmargin=1.4em]
  \item \textbf{Después de un ajuste} --- se genera una nueva simulación con los parámetros ajustados. Las curvas se trazan a partir de los puntos simulados y los datos reales se superponen en rojo.
  \item \textbf{Después de Representar curvas (sim.)} --- las curvas se trazan a partir de la última simulación y, si hay un conjunto de datos reales cargado, estos se superponen en rojo.
  \item \textbf{Cuando está desactivada} --- las curvas utilizan únicamente los datos simulados o ajustados, sin superposición.
\end{enumerate}

\begin{tipbox}[Aplicación]
Esto resulta útil cuando el conjunto de datos experimentales es pequeño: las curvas de percentiles quedan mejor definidas mediante una simulación densa ajustada a los parámetros estimados.
\end{tipbox}

% ════════════════════════════════════════════════════════════
\section{Restablecer la aplicación}
% ════════════════════════════════════════════════════════════

El botón \textbf{Restablecer y cargar archivo} (en \textbf{Ajuste y estimación}) realiza un restablecimiento completo:

\begin{enumerate}[leftmargin=1.4em]
  \item Muestra un cuadro de diálogo de confirmación.
  \item Borra todos los datos, resultados, gráficos y ajustes.
  \item Vuelve a cargar la página.
  \item Abre el cuadro de diálogo para seleccionar un nuevo archivo CSV.
\end{enumerate}

\begin{warnbox}[Irreversible]
Esta acción no se puede deshacer. Se perderá todo el trabajo no guardado, incluidas las figuras que no se hayan exportado.
\end{warnbox}

% ════════════════════════════════════════════════════════════
\section{Interpretar la consola}
% ════════════════════════════════════════════════════════════

La consola (en la parte inferior de la página) utiliza un código de colores:

\begin{center}
\begin{tabular}{ll}
\toprule
\textbf{Color} & \textbf{Significado} \\
\midrule
\textcolor{greenhl}{Verde \textbf{[OK]}} & Operación correcta, finalización. \\
\textcolor{warnhl}{Orange \textbf{[Warning]}} & Advertencia: problema no fatal. \\
\textcolor{redhl}{Rojo \textbf{[Error]}} & Error: la acción ha fallado. \\
\textcolor{bluehl}{Blue} & Mensaje informativo. \\
\textcolor{accent}{\textbf{Gold bold}} & Encabezado de sección / inicio de operación. \\
\bottomrule
\end{tabular}
\end{center}

Mensajes importantes a los que prestar atención:

\begin{itemize}[leftmargin=1.4em]
  \item \textit{``Reference auto-corrected to \ldots''} --- el valor de referencia no era factible y la aplicación lo corrigió.
  \item \textit{``Rejection rate: XX\%''} --- porcentaje de intentos de simulación rechazados por las restricciones físicas.
  \item \textit{``Only XX/YYY valid points''} --- la simulación no pudo alcanzar el número solicitado.
  \item \textit{``Mejor modelo: M\textit{n}''} --- mejor ajuste según el AIC.
\end{itemize}

% ════════════════════════════════════════════════════════════
\section{Solución de problemas}
% ════════════════════════════════════════════════════════════

\begin{description}[leftmargin=0em, style=nextline]
  \item[``No se generaron puntos válidos''] 
  Las restricciones de simulación son demasiado estrictas. Rojouzca $\sM^{\min}$ para acercarlo a $\sMzero$, amplíe los intervalos o compruebe que $\sMzero < \sM^{\min}$.

  \item[``Error al interpretar el CSV''] 
  Compruebe que el encabezado contiene \texttt{sigmam}, \texttt{sigmaM}, \texttt{N}. Compruebe que no haya celdas vacías ni caracteres no numéricos.

  \item[``El ajuste no convergió''] 
  Pruebe con más reinicios (por ejemplo, 10 en la rejilla y 10 en el refinamiento), o introduzca manualmente valores plausibles para $\sMzero$ / $\Nzero$ y utilice \textbf{Ajuste manual}.

  \item[Curves are missing or truncated] 
  El intervalo representado puede ser demasiado estrecho. Amplíe el intervalo de simulación o incorpore más datos.

  \item[``Only 1 valid $\pio$ slice'' in console] 
  Los datos no permiten obtener varias secciones que satisfagan todas las curvas de percentiles en el cuadrante positivo. Amplíe los intervalos de los parámetros $w$ o revise el ajuste.

  \item[La exportación a PDF no hace nada]
  Compruebe que la biblioteca jsPDF esté disponible (acceso de red al CDN o presencia del archivo local). Si no lo está, la consola mostrará un error.
\end{description}

% ════════════════════════════════════════════════════════════
\section{Referencia rápida}
% ════════════════════════════════════════════════════════════

\begin{center}
\begin{tabularx}{\textwidth}{lX}
\toprule
\textbf{Botón} & \textbf{Acción} \\
\midrule
\textgreater Simular datos & Generar puntos sintéticos a partir de los parámetros actuales. \\
\textbf{[Plot]} Representar curvas (sim.) & Dibujar curvas de percentiles a partir de los parámetros de simulación. \\
\textbf{[Plot]} Estime $\sMzero$ \& $\Nzero$ & Búsqueda en rejilla y refinamiento para todos los modelos seleccionados. \\
\textbf{[Fit]} Ajuste manual & Ajustar con los valores de referencia introducidos manualmente. \\
\textbf{[Warning]} Reset and load file & Restablecimiento completo y carga de un nuevo CSV. \\
\textbf{[Download]} CSV / LaTeX & Descargar tablas. \\
\textbf{[Plot]} PDF / \textbf{[Plot]} PNG & Descargar figuras. \\
\bottomrule
\end{tabularx}
\end{center}

\begin{notebox}[Flujo de trabajo recomendado]
\begin{enumerate}
  \item Cargue el archivo CSV.
  \item Estime $\sMzero$ \& $\Nzero$ (opción 1, valores predeterminados).
  \item Examine la pestaña Comparación de modelos e identifique el mejor modelo.
  \item Revise las Curvas de percentiles y ajuste el control de opacidad si es necesario.
  \item Marque \textit{Incluir datos simulados} y vuelva a representar las curvas para obtener una familia más densa.
  \item Exportare las figuras y tablas para elaborar informes.
\end{enumerate}
\end{notebox}

% ════════════════════════════════════════════════════════════
\section*{Historial de versiones}
\addcontentsline{toc}{section}{Historial de versiones}
% ════════════════════════════════════════════════════════════

\textbf{v8} --- Garantía de la rama física en la simulación; informe de la tasa de rechazo; estimación independiente por modelo (opción 1); exportación de PDF por figura; vistas previas de exportación; validación estricta de las secciones de $\pio$-para las curvas de percentiles; control deslizante de opacidad; visualización de la ecuación del modelo.

\end{document}
